\documentclass{tietreport}

% Import images
\usepackage{graphicx}

% Bibliography configuration
\usepackage[
  date=year,
  style=alphabetic,
  backend=biber,
  sorting=nty,
  maxnames=5,
  minnames=3,
  maxalphanames=4,
  minalphanames=3,
  backref=true,
  doi=false,
  isbn=false,
  url=false,
  eprint=false
]{biblatex}
\addbibresource{references.bib}

%% ==================================================
%% PROJECT FIELDS
%% ==================================================

\TitlePageHeader{%
  UCS503P - Software Engineering Project%
}

\ProjectTitle{%
  Faculty Office-Hours Scheduler%
}

\ProjectSubTitle{%
  A Conflict-Free Booking System for Faculty and TA Office Hours%
}

\author{%
  \texttt{1024170227} Manan Garg \\
  \texttt{1024270355} Vishesh \\
  \texttt{1024170240} Aryan \\
  \texttt{1024170244} Dravinam%
}

\TitlePageSubText{%
  Group: \texttt{XXXX} BE Third Year, CSE%
}

\AdvisorName{%
  Dr. Raghav B. Venkataramaiyer%
}

\TitlePageFooterText{{%
    \large Prototype-Stage Evaluation \\[0.2cm]
    \bfseries Computer Science and Engineering
    Department%
  } \\[0.15cm]
  Thapar Institute of Engineering and Technology,
  Patiala%
}

%% ==================================================

\begin{document}

\maketitle
\tableofcontents

% ===== CHAPTER 1: INTRODUCTION =====

\chapter{Introduction}

\section{Background and Context}

Students at Thapar routinely need short, one-on-one time
with professors or teaching assistants to clear doubts on
coursework. In practice, this coordination is entirely
informal: students learn of a faculty member's
availability by word of mouth, wait outside offices
hoping to catch a free moment, or exchange emails and
WhatsApp messages to arrange a meeting. There is no
shared, structured record of who is free, when, or who
has already booked that time.

This informality creates avoidable friction for both
sides. Multiple students may show up expecting the same
informal slot, while other free time on a faculty
member's schedule goes completely unused. Since there is
no record of what was agreed, there is no accountability
for missed sessions -- a student who does not show up
leaves no trace, and a faculty member has no way to
measure how often this happens.

\subsection{Problem Statement}

There is no lightweight, structured system through which
faculty and TAs can publish their available time slots
and have students reserve them without risk of
double-booking or informal miscommunication. This project
addresses that gap directly: it does not attempt to
replace scheduling infrastructure at the institutional
level, but provides a minimal, deployable booking layer
that any faculty member or TA can adopt voluntarily. This
report was prepared using \LaTeX{} \cite{latex2e}.

\section{Project Objectives}

\begin{enumerate}
  \item \textbf{Objective 1:} Allow faculty and TAs to
    publish available office-hour slots through a simple
    web interface, without needing any change to
    institutional systems.
  \item \textbf{Objective 2:} Allow students to browse and
    book an open slot with a single action, with the
    system guaranteeing that no two students can ever book
    the same slot, even if they attempt to do so at the
    same instant.
  \item \textbf{Objective 3:} Track booking outcomes
    (completed or no-show) so that no-show rate, slot
    utilization, and average wait time before a slot is
    available can be measured objectively from logged
    data, rather than estimated.
\end{enumerate}

These objectives are specific (tied to concrete actions:
publish, book, track), measurable (each has a
corresponding numeric metric reported in Chapter 3),
achievable within the tools and skills available to the
team, relevant to a real and recurring problem on campus,
and time-bound to the semester-long scope of this course.

% ===== CHAPTER 2: METHODOLOGY & DESIGN =====

\chapter{Methodology and Design}

\section{System Architecture}

The system follows a standard three-tier web
architecture: a browser-based frontend, a REST API
backend, and a relational database. This separation
allows the booking logic (the part of the system requiring
the most care, see Section~\ref{sec:booking-logic}) to be
implemented and tested independently of the user
interface.

\subsection{High-Level Design}

Three actors interact with the system: \textbf{Student},
\textbf{Faculty/TA}, and \textbf{Admin}. Figure
\ref{fig:usecase} shows the use case diagram capturing
each actor's interactions with the system.

\begin{figure}[h]
\centering
\includegraphics[width=0.85\textwidth]{figures/usecase.png}
\caption{Use case diagram for the Office-Hours Scheduler}
\label{fig:usecase}
\end{figure}

The system's core entities and their relationships are
shown in the class diagram in Figure \ref{fig:class}.
\texttt{User} is a base type from which \texttt{Student},
\texttt{Faculty}, and \texttt{Admin} inherit, reflecting
the role-based access control implemented in the backend.
A \texttt{Faculty} publishes many \texttt{Slot}s; a
\texttt{Student} makes many \texttt{Booking}s; and each
\texttt{Slot} is reserved by at most one \texttt{Booking}
-- this 0..1 cardinality is precisely what the database
uniqueness constraint described in Section
\ref{sec:booking-logic} enforces.

\begin{figure}[h]
\centering
\includegraphics[width=0.8\textwidth]{figures/classdiagram.png}
\caption{Class diagram showing core entities and relationships}
\label{fig:class}
\end{figure}

Figure \ref{fig:activity} traces the control flow through
the booking process, including the availability check that
prevents double-booking.

\begin{figure}[h]
\centering
\includegraphics[width=0.6\textwidth]{figures/activity.png}
\caption{Activity diagram for the slot-booking workflow}
\label{fig:activity}
\end{figure}

\subsection{Technical Stack}

\begin{itemize}
  \item \textbf{Framework/Language:} Node.js with
    Express for the backend REST API; plain HTML/JS for
    the prototype frontend (React planned for the final
    submission).
  \item \textbf{Database:} SQLite for the prototype stage,
    chosen for zero-configuration local development;
    PostgreSQL planned for the deployed final version.
  \item \textbf{Tools:} Node's built-in test runner
    (\texttt{node:test}) for backend unit tests; Git and
    GitHub for version control; GitHub Actions planned for
    CI/CD in the final stage.
  \item \textbf{Libraries:} \texttt{better-sqlite3} for
    synchronous, transactional database access;
    \texttt{jsonwebtoken} for stateless authentication;
    \texttt{bcryptjs} for password hashing.
\end{itemize}

\section{Implementation Details}

\subsection{Module 1: Authentication and Role Management}

Users sign up with a name, email, password, and role
(\texttt{student}, \texttt{faculty}, or \texttt{admin}).
Passwords are hashed with bcrypt before storage. On
successful login, the backend issues a signed JSON Web
Token (JWT) encoding the user's ID, role, and name. Every
subsequent API request carries this token, and a
middleware layer verifies it and enforces role-based
access -- for example, only a user with role
\texttt{faculty} may call the slot-publishing endpoint.

\subsection{Module 2: Conflict-Free Booking Logic}
\label{sec:booking-logic}

This is the most technically significant part of the
system, and the part we tested most rigorously. The risk
is a \textit{race condition}: if two students attempt to
book the same slot at nearly the same instant, a naive
implementation (check availability, then separately write
the booking) could let both succeed, since the check for
student B might run before student A's write completes.

We address this with two mechanisms working together:

\begin{enumerate}
  \item A \texttt{UNIQUE} constraint on the
    \texttt{slot\_id} column of the \texttt{bookings}
    table at the database level, so the database itself
    physically cannot store two bookings for the same
    slot, regardless of application logic.
  \item The check-and-reserve operation (verify the slot
    is open, then mark it booked, then insert the booking
    record) is wrapped in a single atomic database
    transaction, so no other operation can interleave
    between the check and the write.
\end{enumerate}

Listing~\ref{lst:booking} shows the core of this
implementation.

\begin{verbatim}
const bookSlot = db.transaction((slotId, studentId) => {
  const slot = db.prepare(
    "SELECT * FROM slots WHERE id = ? AND status = 'open'"
  ).get(slotId);
  if (!slot) throw new Error('SLOT_UNAVAILABLE');

  db.prepare(
    "UPDATE slots SET status='booked' WHERE id=? AND status='open'"
  ).run(slotId);

  return db.prepare(
    'INSERT INTO bookings (slot_id, student_id) VALUES (?,?)'
  ).run(slotId, studentId);
});
\end{verbatim}
\begin{center}
\small Listing \refstepcounter{figure}\thefigure: Transactional check-and-reserve booking logic \label{lst:booking}
\end{center}

\subsection{Module 3: Slot Publishing and Analytics}

Faculty publish slots by submitting a timestamp; the
system rejects duplicate slots for the same faculty member
at the same time via a second uniqueness constraint on
\texttt{(faculty\_id, slot\_time)}. An analytics endpoint,
restricted to the \texttt{admin} role, aggregates booking
outcomes into the metrics defined in our proposal: total
bookings, no-show rate, and slot utilization rate.

% ===== CHAPTER 3: RESULTS & EVALUATION =====

\chapter{Results and Evaluation}

\section{Testing Methodology}

\begin{itemize}
  \item \textbf{Unit Testing:} The booking transaction
    logic (Section \ref{sec:booking-logic}) is covered by
    automated tests using Node's built-in test runner,
    run against an in-memory SQLite database for speed and
    isolation.
  \item \textbf{Integration Testing:} The full request
    flow -- signup, login, slot publishing, booking, and
    conflict handling -- was verified end-to-end via HTTP
    requests against a running instance of the backend.
  \item \textbf{Manual Testing:} The prototype frontend was
    used manually to sign up as a faculty member, publish
    a slot, sign up as two different students, and attempt
    to book the same slot from both accounts.
\end{itemize}

\section{Performance Metrics}

The central engineering claim of this project is that the
booking system prevents double-booking under concurrent
access. Table \ref{tab:results} reports the outcome of
both the automated test suite and a manual two-student
conflict test performed against the running prototype.

\begin{table}[h]
\centering
\begin{tabular}{|l|l|l|}
  \hline
  \textbf{Test} & \textbf{Expected} & \textbf{Observed} \\
  \hline
  Single booking on open slot & Success (200) & Success (200) \\
  \hline
  Second booking on same slot & Rejected (409) & Rejected (409) \\
  \hline
  Booking non-existent slot & Rejected & Rejected \\
  \hline
  Automated unit tests (2 cases) & 2/2 pass & 2/2 pass \\
  \hline
\end{tabular}
\caption{Booking conflict-prevention test results}
\label{tab:results}
\end{table}

In the manual two-student test, Student~A's booking
request for a published slot returned a confirmed booking.
Student~B's subsequent request for the identical slot was
rejected with the message ``This slot was just booked by
someone else'' and HTTP status 409, confirming that the
uniqueness constraint and transactional logic behave as
designed.

\section{Analysis}

At this prototype stage, we have met Objectives 1 and 2 in
full: faculty can publish slots, and students can book
them with a verified guarantee against double-booking. The
core technical risk identified in our proposal --
concurrent booking conflicts -- has been directly tested
and resolved, rather than assumed to work.

Objective 3 (measurable no-show rate and utilization
tracking) is partially implemented: the analytics endpoint
computes these metrics correctly from logged data, but we
have not yet run a real pilot with actual faculty and
students, so the numbers reported by the endpoint are
currently based on test data rather than genuine usage.
This is the primary gap to close before the final
submission.

The main challenge encountered was deciding how to
simulate concurrent requests in an automated test; we
resolved this by testing the transaction's correctness
directly (a booking attempt on an already-booked slot must
raise an error) rather than attempting to physically
race two HTTP requests against each other, which would
introduce test flakiness without adding confidence beyond
what the transactional guarantee already provides.

% ===== CHAPTER 4: CONCLUSION =====

\chapter{Conclusion}

\section{Summary of Work}

At the prototype stage, we have implemented and tested the
core booking workflow of the Faculty Office-Hours
Scheduler: authentication with role-based access, slot
publishing by faculty, and conflict-free slot booking by
students, backed by a database-level uniqueness constraint
and an atomic transaction. This directly satisfies
Objectives 1 and 2 from Chapter 1 and provides verified
evidence, not just a design claim, that the system's
central technical risk has been addressed.

\section{Key Findings}

\begin{itemize}
  \item \textbf{Finding 1:} A database-level uniqueness
    constraint combined with a single atomic transaction
    is sufficient to prevent double-booking, without
    needing any application-level locking mechanism.
  \item \textbf{Finding 2:} Separating the booking
    transaction logic from the HTTP layer made it possible
    to unit-test the core correctness guarantee in
    isolation, independent of the web framework.
  \item \textbf{Finding 3:} SQLite was sufficient for rapid
    prototype iteration and testing; migrating to
    PostgreSQL for the final, deployed version should be a
    low-risk change since the transactional pattern used
    is standard SQL.
\end{itemize}

\section{Future Work and Recommendations}

\begin{enumerate}
  \item Migrate the database layer from SQLite to
    PostgreSQL for the final submission, to match the
    production-oriented stack described in our proposal.
  \item Rebuild the frontend in React, replacing the
    current minimal HTML/JS prototype, to support a more
    complete faculty availability console and admin
    dashboard.
  \item Add email/in-app reminder notifications ahead of a
    booked session.
  \item Run a real pilot with 3--5 willing faculty/TAs and
    their students to populate the analytics endpoint with
    genuine usage data, and measure actual no-show rate and
    slot utilization against the targets set in our
    proposal.
  \item Set up the GitHub Actions CI/CD pipeline to run the
    existing test suite automatically on every push, and
    deploy the backend to a staging environment (Render or
    Railway).
\end{enumerate}

\section{Final Remarks}

Building and testing the booking logic first, before any
interface polish, let us verify the project's central
engineering claim early and with confidence. The remaining
work is primarily about breadth (a fuller frontend, real
deployment, a genuine pilot) rather than about resolving
open technical risk.

% ===== BIBLIOGRAPHY =====

\printbibliography[heading=bibintoc]

\end{document}
